% ============================================
% Assistant Professor Cover Letter – University of Texas System (Economics)
% ============================================

\documentclass[11pt,letterpaper]{letter}

\usepackage{geometry}
\usepackage{microtype}
\usepackage[hidelinks]{hyperref}

\geometry{left=25mm,right=25mm,top=25mm,bottom=25mm}

\signature{Decory Edwards}
\address{
Department of Economics \\
Johns Hopkins University \\
Baltimore, MD 21218 \\
\texttt{dedwar65@jh.edu}
}

\begin{document}

\begin{letter}{
Search Committee \\
Department of Economics \\
University of Texas System Campus
}

\opening{Dear Members of the Search Committee,}

I am writing to express my strong interest in the tenure-track position in the Department of Economics at the University of Texas System campus. I am a Ph.D.\ candidate in Economics at Johns Hopkins University, advised by Professors Christopher D.~Carroll, Nicholas W.~Papageorge, and Stelios Fourakis, and I expect to receive my degree in May~2026. My research fields are macroeconomics, wealth and income dynamics, and computational economics.

My research agenda focuses on understanding the determinants of wealth inequality and the mechanisms through which heterogeneity in returns to wealth shapes aggregate outcomes. In my job-market paper, I estimate the degree of heterogeneity in individual portfolio returns using micro-data from the Survey of Consumer Finances and simulate a structural model in which households face idiosyncratic return risk. I show that allowing for persistent heterogeneity in returns substantially improves the model’s ability to match the empirical distribution of wealth, offering a tractable way to connect micro-level return dispersion to macroeconomic inequality.

In a second project, I use the Health and Retirement Study to examine the relationship between interpersonal trust and portfolio performance. Building on the literature that finds a hump-shaped relationship between trust and income, I show that a similar relationship holds for individual returns to wealth measured in the HRS data, highlighting a behavioral channel through which social attitudes shape saving and investment outcomes. This work also provides new estimates of the persistent component of returns to net worth, which motivate the calibration of heterogeneity in my structural model.

I am committed to high-quality undergraduate and graduate instruction. I have experience teaching and supporting courses in Principles of Macroeconomics and Intermediate Macroeconomic Theory. My teaching emphasizes economic intuition, careful use of data, and clear connections between theory and policy-relevant questions. I would be well prepared to contribute to the department’s instructional mission and to support students with diverse academic backgrounds in developing quantitative and analytical skills.

I am particularly interested in the University of Texas System’s role as a public institution serving a broad and diverse student population. I would be enthusiastic about contributing to the department’s teaching, research, and service missions and about building a long-term academic career within the University of Texas System.

Thank you very much for your time and consideration. I would welcome the opportunity to discuss my application further.

\closing{Sincerely,}

\end{letter}
\end{document}
